\section{Introduction}

A real homogeneous polynomial $p$ on a finite-dimensional real vector
space is \emph{hyperbolic with respect to} $e$ if $p(e)>0$ and, for every
$x$, all roots of $t\mapsto p(te-x)$ are real. Its closed hyperbolicity
cone is
\[
 K(p,e)=\{x:\text{every root of }t\mapsto p(te-x)
                         \text{ is nonnegative}\}.
\]
G\aa rding proved that this is a convex cone \cite{Garding1959}.
The determinant on real symmetric matrices, with direction the identity,
gives the positive semidefinite cone as a basic example. Hyperbolicity
cones also support logarithmic self-concordant barriers, as shown by
G\"uler \cite{Guler1997}, and hence form a natural setting for convex
optimization beyond semidefinite programming.

A \emph{spectrahedron} is a set defined by a finite affine linear matrix
inequality. A \emph{spectrahedral shadow}, or a set with a finite
semidefinite lift, is a set $S\subset\R^d$ of the form
\begin{equation}\label{eq:lift-definition}
 S=\left\{x\in\R^d:\exists u\in\R^a,
 A_0+\sum_{i=1}^d x_iA_i+\sum_{j=1}^a u_jB_j\succeq0\right\},
\end{equation}
where $a$ and the common size of the real symmetric matrices are finite.
Equality in this definition is exact; taking the closure of the displayed
projection is not part of the definition. The \emph{generalized Lax
conjecture} asserts that every hyperbolicity cone is a spectrahedron.
The weaker \emph{Projected Lax Conjecture}, formulated by Netzer and
Sanyal \cite{NetzerSanyal2015}, asks whether every hyperbolicity cone is
a spectrahedral shadow.

\begin{theorem}\label{thm:main}
There exists a real homogeneous hyperbolic polynomial whose closed
hyperbolicity cone has no finite semidefinite lift.
\end{theorem}

Thus both conjectures have a negative resolution. The assertion excludes
every finite choice in \eqref{eq:lift-definition}, with arbitrary real
coefficients and any number of auxiliary variables. The construction is
an existence argument in high dimension, not a lower bound on the size
of a particular representation.

\subsection{Context and earlier obstructions}

The original Lax question concerned homogeneous polynomials in three
variables \cite{Lax1958}. The determinantal theorem of Helton and
Vinnikov \cite{HeltonVinnikov2007}, together with the equivalence
explained by Lewis, Parrilo, and Ramana \cite{LewisParriloRamana2005},
gave an affirmative answer: a ternary hyperbolic polynomial has a
definite symmetric determinantal representation. The distinction between
a representation of a polynomial and a representation of its cone
becomes essential in higher dimension. Br\"and\'en
\cite{Branden2011} constructed a real-zero polynomial for which no
positive integral power has a definite determinantal representation.
This does not exclude another polynomial with the same positivity
region. Indeed, Kummer \cite{Kummer2016} proved that the related
specialized V\'amos polynomial has a spectrahedral hyperbolicity cone.
Nor does failure of a specified sums-of-squares relaxation, such as
those studied by Kummer, Plaumann, and Vinzant
\cite{KummerPlaumannVinzant2015}, exclude all finite lifts.

Positive representation results cover substantial classes.
Br\"and\'en \cite{Branden2014} proved spectrahedrality for elementary
symmetric hyperbolicity cones. Building on the representability
criteria of Helton and Nie \cite{HeltonNie2010}, Netzer and Sanyal
\cite[Theorem~1.1]{NetzerSanyal2015} proved that a hyperbolicity cone
is a spectrahedral shadow when each nonzero boundary point is a smooth
point of its defining hyperbolic polynomial. Scheiderer
\cite{Scheiderer2025} subsequently obtained second-order-cone lifts
for hyperbolicity cones with Nash-smooth boundary. Recent preprints
prove spectrahedrality for
five-variable hyperbolic cubics \cite{Netzer2026} and for strictly
hyperbolic polynomials \cite{KummerNetzer2026}. These results concern
special families or impose regularity hypotheses; they do not yield
an exact finite lift for every hyperbolicity cone.

For general convex semialgebraic sets, Scheiderer
\cite{Scheiderer2018} established nonexistence of semidefinite lifts
by local positivity obstructions. His finite-dimensional
sum-of-squares pullback criterion and local arguments explain how a
putative lift constrains functions near a smooth point. He explicitly
asked whether these methods could settle the hyperbolicity-cone
question \cite[Section~5.4]{Scheiderer2018}. The argument here is
closely related to that approach, but tests analytic Gram identities
directly by a finite positive moment functional. Moment and
sum-of-squares semidefinite methods are classical
\cite{Lasserre2001}; our functional need not be integration against a
measure. The general relationship between cone lifts and positive
factorizations is developed by Gouveia, Parrilo, and Thomas
\cite{GouveiaParriloThomas2013}. We do not need to invoke their
factorization theorem.

\subsection{The construction and proof strategy}

We begin with a matrix-valued homogeneous quadratic map
$Q:\R^m\to\Sym_r$ that is positive semidefinite at every real input.
Here $\Sym_r$ denotes the real symmetric $r$-by-$r$ matrices.
Section~\ref{sec:cone} associates to $Q$ a homogeneous polynomial
$p_Q$ and proves its hyperbolicity by a symmetric determinant identity
on each line in the input space. It also identifies an exact affine
section and projection of its closed cone:
\[
 E_Q=\{(A,t,y): A\succeq0,\ \ran Q(y)\subseteq\ran A,
                    \ t\ge\tr(A^\dagger Q(y))\}.
\]
The symbol $A^\dagger$ denotes the Moore--Penrose inverse. It is enough
to choose $Q$ so that $E_Q$ has no finite lift.

The difficulty is to force such an obstruction while keeping $Q(y)$
positive semidefinite for \emph{every} $y$. Sections~\ref{sec:hankel}
and~\ref{sec:separation} address this constraint. Let $H$ be the
36-dimensional space of quadratic forms in eight variables. Each
linear functional $y$ on quartic forms determines a Hankel matrix
$Y(y)$ on $H$ by $Y(y)(f,g)=y(fg)$. We construct a rank-twenty
positive semidefinite Hankel matrix whose kernel has strong
multiplication properties. After normalizing translations and dilations,
its geometry yields a twenty-dimensional subspace $W\subset H$ and a
finite set of projective classes of nonzero symmetric forms on $W$.
The class of every nonzero limit of inverses of compressions to $W$
arising in the separation argument belongs to this set. The seed matrix
is established by the finite-field certificate in
Appendix~\ref{app:certificate}; the passage from its rank computations
to a real positive semidefinite matrix is part of the proof.

For a sufficiently large integer $N$, finiteness lets us choose a tuple
$w=(w_1,\ldots,w_N)\in W^N$ with $\sum_i w_i^TCw_i\ne0$ for every
representative $C$ of one of these classes. We then choose
$U\subset w^\perp\subset H^N$ so that convergence to $w$ of vectors
$(I_N\otimes Y(y))u$, $u\in U$, would force one of these pairings
to vanish. Thus $w$ lies outside the closure of those Hankel images.
A symmetric form $g$ can be chosen nonnegative on every such image
and negative at $w$. Using a basis matrix also denoted by $U$, we set
\[
 G(y)=I_N\otimes Y(y),\qquad Q(y)=U^\top G(y)gG(y)U.
\]
Nonnegativity on the Hankel images proves that $Q$ is admissible.
The construction uses $m=330$ and $r=20N-1$, where one may take
$N>10\cdot36^2$.

Section~\ref{sec:obstruction} turns the negative direction of $g$
into a local obstruction. Evaluation functionals at points $(1,s)$,
$s\in\R^7$, give a degree-four polynomial patch $x(s)$ in $E_Q$.
A positive definite moment functional $\Lambda$, applied to
$x(c+\eps s)$, produces points violating the trace inequality of
$E_Q$ by a negative term of order $\eps^4$. At the same time their
$A$ coordinates are bounded below by a positive multiple of
$\eps^4 I$.

If a finite lift existed, semialgebraic selection would provide analytic
lift coordinates and analytic Gram factors on an open part of this
patch. Applying $\Lambda$ to their Taylor polynomials almost preserves
the lift inequality. A mixing argument on the common support of the
feasible pencil values repairs the small error without a strict
feasibility assumption. It yields feasible points within
$O(\eps^{16})$ of the tested patch, too close to remove its order-four
violation. The finite positive extension and finite-jet approximation
lemmas isolate reusable parts of this obstruction. Together with the
finite inverse-limit construction, they supply the proof of
Theorem~\ref{thm:main}.
